% Generated draft: requires mathematical and source review.
\documentclass[11pt]{article}
\usepackage[margin=25mm]{geometry}
\usepackage{fontspec}
\setmainfont{Latin Modern Roman}
\usepackage{amsmath,amssymb,mathtools,unicode-math}
\setmathfont{Latin Modern Math}
\usepackage{enumitem,needspace,xurl,hyperref}
\setlength{\parindent}{0pt}
\setlength{\parskip}{5pt}
\setlength{\emergencystretch}{3em}
\newcommand{\N}{\mathbb{N}}
\newcommand{\Z}{\mathbb{Z}}
\newcommand{\Q}{\mathbb{Q}}
\newcommand{\R}{\mathbb{R}}
\newcommand{\C}{\mathbb{C}}
\newcommand{\F}{\mathbb{F}}
\newcommand{\A}{\mathbb{A}}
\newcommand{\PP}{\mathbb{P}}
\newcommand{\E}{\mathbb{E}}
\newcommand{\eps}{\varepsilon}
\newcommand{\dd}{\,\mathrm{d}}
\begin{document}
\section*{UnsolvedMath ID: 30002052}
\textbf{Source problem code:} catalogue entry
\subsection{Definitions and mathematical statement}
Let \(X\) be a smooth complex algebraic variety; no properness assumption is supplied. Let \(\operatorname{Perf}(X)\) denote the differential graded category of perfect complexes on \(X\), namely complexes locally quasi-isomorphic to bounded complexes of finite-rank locally free \(\mathcal O_X\)-modules. Its algebraic de Rham cohomology is
\[
H^n_{\mathrm{dR}}(X)=\mathbb H^n(X,\Omega^\bullet_{X/\C}),\qquad n\geq 0.
\]
The mixed Hodge structure in question consists of the rational cohomology \(H^n(X^{\mathrm{an}},\Q)\), an increasing weight filtration \(W_\bullet\), and a decreasing Hodge filtration \(F^\bullet\) on its complexification, identified with \(H^n_{\mathrm{dR}}(X)\) by the de Rham comparison. The induced filtration on each \(\operatorname{Gr}^W_m\) defines a pure Hodge structure of weight \(m\).

\textbf{Source question.} For every smooth complex algebraic variety \(X\), can any part of this mixed Hodge structure be recovered from \(\operatorname{Perf}(X)\)? More generally, does the periodic cyclic homology \(HP_\bullet(\mathcal A)\) of every differential graded category \(\mathcal A\) carry a functorial Hodge-like structure? Here periodic cyclic homology means the homology of the periodic cyclic complex associated to the Hochschild mixed complex of \(\mathcal A\).

\textbf{Unresolved formulation.} The source does not specify which part of the mixed Hodge structure is to be recovered, what equivalence relation on differential graded categories the recovery must respect, or whether recovery includes the rational structure, weight filtration, Hodge filtration, or cohomological grading. Consequently, this question does not specify a single precise recovery assertion. Likewise, the general question supplies no base field, size hypotheses, definition of a Hodge-like structure, class of morphisms for functoriality, or compatibility requirements with the geometric case. These choices must be stated before the general question becomes a precise existence conjecture. In particular, the complex base field of the geometric question cannot silently be imposed on its universal formulation.

The original wording allows an algebraic variety \(X\) without explicitly imposing smoothness or a complex base field, whereas the supplied clean statement imposes both. The geometric question above follows the clean statement; the intended scope of the broader original wording remains unresolved.
\subsection{Short English statement}
The problem asks what mixed Hodge information a smooth complex variety's category of perfect complexes determines, and whether periodic cyclic homology of every differential graded category admits functorial Hodge-like data. The recoverable information and the required general structure are unspecified, so this is a research question requiring further formulation and human mathematical review.
\subsection{Sources}
\begin{itemize}
\item Supplied local UnsolvedMath registry record, permanent ID \(30002052\), title \emph{Hodge Structures on Periodic Cyclic Homology}, source problem code \texttt{catalogue entry}, in the Oberwolfach Reports Open Problems collection. No individual proposer is supplied; the proposed year is recorded as 2012.
\item Supplied report citation: \emph{Singularity Theory and Integrable Systems (2013)}, Oberwolfach Reports, DOI \texttt{10.4171/owr/2012/23}, pp.\ 1361--1403. \url{https://doi.org/10.4171/owr/2012/23}. The background attributes the problem to workshop open-problem discussions documented in this report. The citation's 2013 designation and its 2012 DOI/proposed-year information are retained without reconciliation.
\item Shklyarov, \emph{Non-commutative Hodge structures}, as labeled in the supplied literature sources: \url{https://arxiv.org/abs/1107.3156}.
\item Sheridan, \emph{Formulae in noncommutative Hodge theory}, as labeled in the supplied literature sources: \url{https://arxiv.org/abs/1510.03795}.
\item Literature assessment dated exactly \texttt{2026-08-21}: the supplied curation labels the problem partially solved and states that noncommutative Hodge structures and Hodge-to-de Rham type results provide functorial Hodge-like structures for broad classes of differential graded categories, while general canonical recovery of the mixed Hodge structure of every \(\operatorname{Perf}(X)\) is not established. This is supplied dated curation derived from the OWR record, not an independently verified assessment; its time-sensitive status and cited claims require independent verification.
\item The registry labels the statement \texttt{corrected\_verified} and reports comparison with the official Oberwolfach report. This verification claim is supplied metadata and has not been independently checked here.
\end{itemize}
\end{document}
